\begin{appendices}

\renewcommand{\theFancyVerbLine}{%
	\textcolor{red}{\small
		\arabic{FancyVerbLine}}}

\chapter{{\tt RCTTypeBulletproof2} 交易结构}
\label{appendix:RCTTypeBulletproof2}

我们在本附录中展示一笔真实 Monero 交易（类型 {\tt RCTTypeBulletproof2}）的转储，
并附有相关字段的解释说明。

该转储来自区块浏览器 \url{https://xmrchain.net}。也可以通过在非分离模式下运行的 {\tt monerod} 守护进程中执行命令 {\tt print\_tx <TransactionID> +hex +json} 来获取。{\tt <TransactionID>} 是交易的 hash（第 \ref{subsec:transaction-id} 节）。打印的第一行显示了要执行的实际命令。

要复现我们的结果，读者可以执行以下操作：
\begin{enumerate}
    \item 我们需要 Monero 命令行工具（CLI），可在 \url{https://web.getmonero.org/downloads/}（以及其他地方）获取。获取适合你操作系统的"Command Line Tools Only"，将文件移到合适的位置，然后解压。
    \item 打开终端/命令行，导航到解压后创建的文件夹。
    \item 使用 {\tt ./monerod} 运行 {\tt monerod}。Monero blockchain 将开始下载。遗憾的是，目前没有办法在不下载 blockchain 的情况下在自己的系统上打印交易（即不使用区块浏览器）。
    \item 同步完成后，{\tt print\_tx} 等命令就可以使用了。使用 {\tt help} 了解其他命令。
\end{enumerate}

组件 {\tt rctsig\_prunable}，顾名思义，是可以从 blockchain 中 {\sl 修剪} 的。也就是说，一旦区块达成共识且当前链长度排除了所有双花攻击的可能性，就可以修剪整个字段并用其 hash 替换 Merkle 树。

key image 单独存储在交易的不可修剪区域中。这些对于检测双花攻击至关重要，不能被修剪。\\

我们的示例交易有 2 个 input 和 2 个 output，于 2020-03-02 19:01:10 UTC 的时间戳被添加到 blockchain（由其区块的矿工报告）。

\begin{Verbatim}[commandchars=\\\{\}, numbers=left]
print_tx 84799c2fc4c18188102041a74cef79486181df96478b717e8703512c7f7f3349
Found in blockchain at height 2045821
\{
  "version": 2, 
  "unlock_time": 0, 
  "vin": [ \{
      "key": \{
        "amount": 0, 
        "key_offsets": [ 14401866, 142824, 615514, 18703, 5949, 22840, 5572, 16439,
        983, 4050, 320
        ], 
        "k_image": "c439b9f0da76ca0bb17920ca1f1f3f1d216090751752b091bef9006918cb3db4"
      \}
    \}, \{
      "key": \{
        "amount": 0, 
        "key_offsets": [ 14515357, 640505, 8794, 1246, 20300, 18577, 17108, 9824, 581,
        637, 1023
        ], 
        "k_image": "03750c4b23e5be486e62608443151fa63992236910c41fa0c4a0a938bc6f5a37"
      \}
    \}
  ], 
  "vout": [ \{
      "amount": 0, 
      "target": \{
        "key": "d890ba9ebfa1b44d0bd945126ad29a29d8975e7247189e5076c19fa7e3a8cb00"
      \}
    \}, \{
      "amount": 0, 
      "target": \{
        "key": "dbec330f8a67124860a9bfb86b66db18854986bd540e710365ad6079c8a1c7b0"
      \}
    \}
  ], 
  "extra": [ 1, 3, 39, 58, 185, 169, 82, 229, 226, 22, 101, 230, 254, 20, 143,
  37, 139, 28, 114, 77, 160, 229, 250, 107, 73, 105, 64, 208, 154, 182, 158, 200,
  73, 2, 9, 1, 12, 76, 161, 40, 250, 50, 135, 231
  ], 
  "rct_signatures": \{
    "type": 4, 
    "txnFee": 32460000, 
    "ecdhInfo": [ \{
        "amount": "171f967524e29632"
      \}, \{
        "amount": "5c2a1a9f54ccf40b"
      \}], 
    "outPk": [ "fed8aded6914f789b63c37f9d2eb5ee77149e1aa4700a482aea53f82177b3b41",
    "670e086e40511a279e0e4be89c9417b4767251c5a68b4fc3deb80fdae7269c17"]
  \}, 
  "rctsig_prunable": \{
    "nbp": 1, 
    "bp": [ \{
        "A": "98e5f23484e97bb5b2d453505db79caadf20dc2b69dd3f2b3dbf2a53ca280216", 
        "S": "b791d4bc6a4d71de5a79673ed4a5487a184122321ede0b7341bc3fdc0915a796", 
        "T1": "5d58cfa9b69ecdb2375647729e34e24ce5eb996b5275aa93f9871259f3a1aecd", 
        "T2": "1101994fea209b71a2aa25586e429c4c0f440067e2b197469aa1a9a1512f84b7", 
        "taux": "b0ad39da006404ccacee7f6d4658cf17e0f42419c284bdca03c0250303706c03", 
        "mu": "cacd7ca5404afa28e7c39918d9f80b7fe5e572a92a10696186d029b4923fa200", 
        "L": [ "d06404fc35a60c6c47a04e2e43435cb030267134847f7a49831a61f82307fc32",
        "c9a5932468839ee0cda1aa2815f156746d4dce79dab3013f4c9946fce6b69eff",
        "efdae043dcedb79512581480d80871c51e063fe9b7a5451829f7a7824bcc5a0b",
        "56fd2e74ac6e1766cfd56c8303a90c68165a6b0855fae1d5b51a2e035f333a1d",
        "81736ed768f57e7f8d440b4b18228d348dce1eca68f969e75fab458f44174c99",
        "695711950e076f54cf24ad4408d309c1873d0f4bf40c449ef28d577ba74dd86d",
        "4dc3147619a6c9401fec004652df290800069b776fe31b3c5cf98f64eb13ef2c"
        ], 
        "R": [ "7650b8da45c705496c26136b4c1104a8da601ea761df8bba07f1249495d8f1ce",
        "e87789fbe99a44554871fcf811723ee350cba40276ad5f1696a62d91a4363fa6",
        "ebf663fe9bb580f0154d52ce2a6dae544e7f6fb2d3808531b0b0749f5152ddbf",
        "5a4152682a1e812b196a265a6ba02e3647a6bd456b7987adff288c5b0b556ec6",
        "dc0dcb2e696e11e4b62c20b6bfcb6182290748c5de254d64bf7f9e3c38fb46c9",
        "101e2271ced03b229b88228d74b36088b40c88f26db8b1f9935b85fb3ab96043",
        "b14aae1d35c9b176ac526c23f31b044559da75cf95bc640d1005bfcc6367040b"
        ], 
        "a": "4809857de0bd6becdb64b85e9dfbf6085743a8496006b72ceb81e01080965003", 
        "b": "791d8dc3a4ddde5ba2416546127eb194918839ced3dea7399f9c36a17f36150e", 
        "t": "aace86a7a1cbdec3691859fa07fdc83eed9ca84b8a064ca3f0149e7d6b721c03"
      \}
    ], 
    "MGs": [ \{
        "ss": [[ "d7a9b87cfa74ad5322eedd1bff4c4dca08bcff6f8578a29a8bc4ad6789dee106",
        "f08e5dfade29d2e60e981cb561d749ea96ddc7e6855f76f9b807842d1a17fe00"],
        ["de0a86d12be2426f605a5183446e3323275fe744f52fb439041ad2d59136ea0b",
        "0028f97976630406e12c54094cbbe23a23fe5098f43bcae37339bfc0c4c74c07"],
        ["f6eef1f99e605372cb1ec2b3dd4c6e56a550fec071c8b1d830b9fda34de5cc05",
        "cd98fc987374a0ac993edf4c9af0a6f2d5b054f2af601b612ea118f405303306"],
        ["5a8437575dae7e2183a1c620efbce655f3d6dc31e64c96276f04976243461e08",
        "5090103f7f73a33024fbda999cd841b99b87c45fa32c4097cdc222fa3d7e9502"],
        ["88d34246afbccbd24d2af2ba29d835813634e619912ea4ca194a32281ac14e0e",
        "eacdf59478f132dd8dbb9580546f96de194092558ffceeff410ee9eb30ce570e"],
        ["571dab8557921bbae30bda9b7e613c8a0cff378d1ec6413f59e4972f30f2470d",
        "5ca78da9a129619299304d9b03186233370023debfdaddcd49c1a338c1f0c50d"],
        ["ac8dbe6bb28839cf98f02908bd1451742a10c713fdd51319f2d42a58bf1d7507",
        "7347bf16cba5ee6a6f2d4f6a59d1ed0c1a43060c3a235531e7f1a75cd8c8530d"],
        ["b8876bd3a5766150f0fbc675ba9c774c2851c04afc4de0b17d3ac4b6de617402",
        "e39f1d2452d76521cbf02b85a6b626eeb5994f6f28ce5cf81adc0ff2b8adb907"],
        ["1309f8ead30b7be8d0c5932743b343ef6c0001cef3a4101eae98ffde53f46300",
        "370693fa86838984e9a7232bca42fd3d6c0c2119d44471d61eee5233ba53c20f"],
        ["80bc2da5fc5951f2c7406fce37a7aa72ffef9cfa21595b1b68dfab4b7b9f9f0c",
        "c37137898234f00bce00746b131790f3223f97960eefe67231eb001092f5510c"],
        ["01c89e07571fd365cac6744b34f1b44e06c1c31cbf3ee4156d08309345fdb20e",
        "a35c8786695a86c0a4e677b102197a11dadc7171dd8c2e1de90d828f050ec00f"]], 
        "cc": "0d8b70c600c67714f3e9a0480f1ffc7477023c793752c1152d5df0813f75ff0f"
      \}, \{
        "ss": [[ "4536e585af58688b69d932ef3436947a13d2908755d1c644ca9d6a978f0f0206",
        "9aab6509f4650482529219a805ee09cd96bb439ee1766ced5d3877bf1518370b"],
        ["5849d6bf0f850fcee7acbef74bd7f02f77ecfaaa16a872f52479ebd27339760f",
        "96a9ec61486b04201313ac8687eaf281af59af9fd10cf450cb26e9dc8f1ce804"],
        ["7fe5dcc4d3eff02fca4fb4fa0a7299d212cd8cd43ec922d536f21f92c8f93f00",
        "d306a62831b49700ae9daad44fcd00c541b959da32c4049d5bdd49be28d96701"],
        ["2edb125a5670d30f6820c01b04b93dd8ff11f4d82d78e2379fe29d7a68d9c103",
        "753ac25628c0dada7779c6f3f13980dfc5b7518fb5855fd0e7274e3075a3410c"],
        ["264de632d9cb867e052f95007dfdf5a199975136c907f1d6ad73061938f49c01",
        "dd7eb6028d0695411f647058f75c42c67660f10e265c83d024c4199bed073d01"],
        ["b2ac07539336954f2e9b9cba298d4e1faa98e13e7039f7ae4234ac801641340f",
        "69e130422516b82b456927b64fe02732a3f12b5ee00c7786fe2a381325bf3004"],
        ["49ea699ca8cf2656d69020492cdfa69815fb69145e8f922bb32e358c23cebb0f",
        "c5706f903c04c7bed9c74844f8e24521b01bc07b8dbf597621cceeeb3afc1d0c"],
        ["a1faf85aa942ba30b9f0511141fcab3218c00953d046680d36e09c35c04be905",
        "7b6b1b6fb23e0ee5ea43c2498ea60f4fcf62f70c7e0e905eb4d9afa1d0a18800"],
        ["785d0993a70f1c2f0ac33c1f7632d64e34dd730d1d8a2fb0606f5770ed633506",
        "e12777c49ffc3f6c35d27a9ccb3d9b8fed7f0864a880f7bae7399e334207280e"],
        ["ab31972bf1d2f904d6b0bf18f4664fa2b16a1fb2644cd4e6278b63ade87b6d09",
        "1efb04fe9a75c01a0fe291d0ae00c716e18c64199c1716a086dd6e32f63e0a07"],
        ["a6f4e21a27bf8d28fc81c873f63f8d78e017666adbf038da0b83c2ad04ef6805",
        "c02103455f93c2d7ec4b7152db7de00d1c9e806b1945426b6773026b4a85dd03"]], 
        "cc": "d5ac037bb78db41cf924af713b7379c39a4e13901d3eac017238550a1a3b910a"
      \}],
    "pseudoOuts": [ "b313c1ae9ca06213684fbdefa9412f4966ad192bc0b2f74ed1731381adb7ab58",
    "7148e7ef5cfd156c62a6e285e5712f8ef123575499ff9a11f838289870522423"]
  \}
\}
\end{Verbatim}



\section*{交易组件}
	
\begin{itemize}
    \item （第 2 行）- {\tt print\_tx} 命令会报告找到交易的区块，此处复制以供演示。
	\item {\tt version}（第 4 行）- 交易格式/时代版本；`2' 对应 RingCT。
	\item {\tt unlock\_time}（第 5 行）- 阻止交易的 output 在指定时间到达前被花费。它要么是区块高度，要么（如果数字大于 UNIX 时间起始点）是 UNIX 时间戳。默认为零，表示没有限制。
	\item {\tt vin}（第 6-23 行）- input 列表（此处有两个）
	\item {\tt amount}（第 8 行）- 已弃用的（旧版）金额字段，用于类型 1 的交易
	\item {\tt key\_offset}（第 9 行）- 这允许验证者在 blockchain 中查找环成员的 key 和承诺，并证明这些成员是合法的。第一个偏移量是 blockchain 历史中的绝对位置，随后的每个偏移量相对于前一个。例如，真实偏移量为 \{7,11,15,20\} 时，blockchain 记录为 \{7,4,4,5\}。验证者通过 (7+4+4+5 = 20) 计算最后一个偏移量（第 \ref{subsec:space-and-ver-rcttypefull} 节）。
	\item {\tt k\_image}（第 12 行）- 来自第 \ref{sec:MLSAG} 节的 key image \(\tilde{K_j}\)，其中 $j = 1$，因为这是第一个 input。
	\item {\tt vout}（第 24-35 行）- output 列表（此处有两个）
	\item {\tt amount}（第 25 行）- 已弃用的金额字段，用于类型 1 的交易
	\item {\tt key}（第 27 行）- 第 \ref{sec:one-time-addresses} 节所述的 output $t = 0$ 的一次性目标 key
	\item {\tt extra}（第 36-39 行）- 其他\marginnote{src/crypto- note\_basic/ tx\_extra.h} 数据，包括 {\em 交易公钥}（即第 \ref{sec:one-time-addresses} 节的共享秘密 $r G$）以及第 \ref{sec:integrated-addresses} 节中的加密付款 ID。通常的编码方式如下：每个数字是一个字节（取值范围 0-255），该字段中每种类型的数据都有一个"标签"和"长度"。标签指示接下来是什么信息，长度指示该信息占用多少字节。第一个数字总是标签。这里 `1' 表示"交易公钥"。交易公钥始终为 32 字节，所以不需要包含长度。32 个数字之后我们找到新标签 `2'，表示"extra nonce"，其长度为 `9'，下一个字节 `1' 表示一个 8 字节的加密付款 ID（出于某种原因，extra nonce 内部可以有子字段）。再过 8 个字节，就到了这个 extra 字段的末尾。更多详情请参见 \cite{extra-field-stackexchange}。（注意：在原始 Cryptonote 规范中，第一个字节表示字段的大小。Monero 不使用这种方式。）\cite{tx-extra-field}
	\item {\tt rct\_signatures}（第 40-50 行）- 签名数据的第一部分
	\item {\tt type}（第 41 行）- 签名类型；{\tt RCTTypeBulletproof2} 是类型 4。已弃用的 RingCT 类型 {\tt RCTTypeFull} 和 {\tt RCTTypeSimple} 分别是 1 和 2。矿工交易使用 {\tt RCTTypeNull}，类型 0。
	\item {\tt txnFee}（第 42 行）- 以明文显示的交易费用，此例中为 0.00003246 XMR
	\item {\tt ecdhInfo}（第 43-47 行）- "椭圆曲线 Diffie-Hellman 信息"：每个 output $t \in \{0, ..., p-1\}$ 的混淆金额；此处 $p = 2$
    \item {\tt amount}（第 44 行）- 第 \ref{sec:pedersen_monero} 节所述的 $t = 0$ 处的 {\sl amount} 字段
    \item {\tt outPk}（第 48-49 行）- 每个 output 的承诺，见第 \ref{sec:ringct-introduction} 节

    \item {\tt rctsig\_prunable}（第 51-132 行）- 签名数据的第二部分
    \item {\tt nbp}（第 52 行）- 此交易中 Bulletproof 范围证明的数量
    \item {\tt bp}（第 53-80 行）- Bulletproof 证明元素（本文未探讨 Bulletproofs，因此不作进一步枚举）\vspace{.175cm}
    \[\Pi_{BP} = (A, S, T_1, T_2, \tau_x, \mu, \mathbb{L}, \mathbb{R}, a, b, t)\]
    \item {\tt MGs}（第 81-129 行）- MLSAG 签名
    \item {\tt ss}（第 82-103 行）- 第一个 input 的 MLSAG 签名中的分量 \(r_{i,1}\) 和 \(r_{i,2}\)\vspace{.175cm}
    \[\sigma_j(\mathfrak{m}) = (c_1, r_{1, 1}, r_{1, 2}, ..., r_{v+1, 1}, r_{v+1, 2})\]
    \item {\tt cc}（第 104 行）- 上述 MLSAG 签名中的分量 \(c_1\)
    \item {\tt pseudoOuts}（第 130-131 行）- 伪 output 承诺 $C'^a_j$，如第 \ref{sec:pedersen_monero} 节所述。请记住，这些承诺之和将等于本笔交易的两个 output 承诺之和（加上交易费用承诺 $f H$）。
\end{itemize}




\chapter{区块内容}
\label{appendix:block-content}

在本附录中，我们展示一个示例区块的结构，即创世区块之后的第 1582196 个区块。该区块有 5 笔交易，于 2018-05-27 21:56:01 UTC 的时间戳被添加到 blockchain（由其区块的矿工报告）。

\begin{Verbatim}[commandchars=\\\{\}, numbers=left]
print_block 1582196
timestamp: 1527458161
previous hash: 30bb9b475a08f2ea6fe07a1fd591ea209a7f633a400b2673b8835a975348b0eb
nonce: 2147489363
is orphan: 0
height: 1582196
depth: 2
hash: 50c8e5e51453c2ab85ef99d817e166540b40ef5fd2ed15ebc863091ca2a04594
difficulty: 51333809600
reward: 4634817937431
\{
  "major_version": 7,
  "minor_version": 7,
  "timestamp": 1527458161,
  "prev_id": "30bb9b475a08f2ea6fe07a1fd591ea209a7f633a400b2673b8835a975348b0eb",
  "nonce": 2147489363,
  "miner_tx": \{
    "version": 2,
    "unlock_time": 1582256,
    "vin": [ \{
        "gen": \{
          "height": 1582196
        \}
      \}
    ],
    "vout": [ \{
        "amount": 4634817937431,
        "target": \{
          "key": "39abd5f1c13dc6644d1c43db68691996bb3cd4a8619a37a227667cf2bf055401"
        \}
      \}
    ],
    "extra": [ 1, 89, 148, 148, 232, 110, 49, 77, 175, 158, 102, 45, 72, 201, 193,
    18, 142, 202, 224, 47, 73, 31, 207, 236, 251, 94, 179, 190, 71, 72, 251, 110, 
    134, 2, 8, 1, 242, 62, 180, 82, 253, 252, 0
    ],
    "rct_signatures": \{
      "type": 0
    \}
  \},
  "tx_hashes": [ "e9620db41b6b4e9ee675f7bfdeb9b9774b92aca0c53219247b8f8c7aecf773ae",
                 "6d31593cd5680b849390f09d7ae70527653abb67d8e7fdca9e0154e5712591bf",
                 "329e9c0caf6c32b0b7bf60d1c03655156bf33c0b09b6a39889c2ff9a24e94a54",
                 "447c77a67adeb5dbf402183bc79201d6d6c2f65841ce95cf03621da5a6bffefc",
                 "90a698b0db89bbb0704a4ffa4179dc149f8f8d01269a85f46ccd7f0007167ee4"
  ]
\}
\end{Verbatim}



\section*{区块组件}

\begin{itemize}
	\item （第 2-10 行）- 由软件收集的区块信息，严格来说并不是区块本身的组成部分。
    \item {\tt is orphan}（第 5 行）- 表示该区块是否为孤块。节点在分叉情况下通常会存储所有分支，当累计难度赢家出现时丢弃不必要的分支，从而使这些区块成为孤块。
    \item {\tt depth}（第 7 行）- 在 blockchain 副本中，任何给定区块的 depth 是其相对于最新区块的链上深度。
    \item {\tt hash}（第 8 行）- 该区块的区块 ID。
    \item {\tt difficulty}（第 9 行）- 难度不存储在区块中，因为用户可以从区块时间戳和第 \ref{sec:difficulty} 节中的规则计算出 {\em 所有} 区块的难度。
    \item {\tt major\_version}（第 12 行）- 对应于验证此区块所使用的协议版本。
    \item {\tt minor\_version}（第 13 行）- 最初旨在作为矿工之间的投票机制，现在只是重复 {\tt major\_version}。由于该字段占用空间不大，开发者可能认为弃用它不值得花精力。
    \item {\tt timestamp}（第 14 行）- 该区块 UTC 时间戳的整数表示，由区块的矿工报告。
    \item {\tt prev\_id}（第 15 行）- 前一个区块的区块 ID。Monero blockchain 的精髓就在于此。
    \item {\tt nonce}（第 16 行）- 该区块矿工用来通过难度目标的 nonce。任何人都可以重新计算工作量证明并验证 nonce 是否有效。
    \item {\tt miner\_tx}（第 17-40 行）- 该区块的矿工交易。
    \item {\tt version}（第 18 行）- 交易格式/时代版本；`2' 对应 RingCT。
    \item {\tt unlock\_time}（第 19 行）- 矿工交易的 output 直到第 1582256 个区块才能被花费，即再开采 59 个区块之后（这是 60 个区块的锁定时间，因为只有在 60 个区块时间间隔过去后才能花费，例如 $2*60 = 120$ 分钟）。
    \item {\tt vin}（第 20-25 行）- 矿工交易的 input。没有 input，因为矿工交易用于生成区块奖励并收集交易费用。
    \item {\tt gen}（第 21 行）- "generate"的缩写。
    \item {\tt height}（第 22 行）- 生成该矿工交易区块奖励的区块高度。每个区块高度只能生成一次区块奖励。
    \item {\tt vout}（第 26-32 行）- 矿工交易的 output。
    \item {\tt amount}（第 27 行）- 矿工交易分配的金额，包含区块奖励和本区块交易的费用。以原子单位记录。
    \item {\tt key}（第 29 行）- 分配矿工交易金额所有权的一次性地址。
    \item {\tt extra}（第 33-36 行）- 矿工交易的额外信息，包括交易公钥。
    \item {\tt type}（第 38 行）- 交易类型，此处 `0' 表示 {\tt RCTTypeNull}，即矿工交易。
    \item {\tt tx\_hashes}（第 41-46 行）- 本区块包含的所有交易 ID（不包括矿工交易 ID，其为 {\tt 06fb3e1cf889bb972774a8535208d98db164394ef2b14ecfe74814170557e6e9}）。
\end{itemize}




\chapter{创世区块 (Genesis Block)}
\label{appendix:genesis-block}

在本附录中，我们展示 Monero 创世区块的结构。该区块有 0 笔交易（它只是将第一个区块奖励发送给 thankful\_for\_today \cite{bitmonero-launched}）。Monero 的创始人没有添加时间戳，这可能是 Bytecoin（Monero 代码由此分叉而来）的遗留痕迹，其创建者显然试图隐藏大量的预挖 \cite{monero-history}，并且可能运营着一个与加密货币相关的软件和服务的可疑网络 \cite{bytecoin-network}。

第 1 个区块于 2014-04-18 10:49:53 UTC 的时间戳被添加到 blockchain（由其区块的矿工报告），因此我们可以假设创世区块是在同一天创建的。这与 thankful\_for\_today 宣布的发布日期一致 \cite{bitmonero-launched}。

\begin{Verbatim}[commandchars=\\\{\}, numbers=left]
print_block 0
timestamp: 0
previous hash: 0000000000000000000000000000000000000000000000000000000000000000
nonce: 10000
is orphan: 0
height: 0
depth: 1580975
hash: 418015bb9ae982a1975da7d79277c2705727a56894ba0fb246adaabb1f4632e3
difficulty: 1
reward: 17592186044415
\{
  "major_version": 1,
  "minor_version": 0,
  "timestamp": 0,
  "prev_id": "0000000000000000000000000000000000000000000000000000000000000000",
  "nonce": 10000,
  "miner_tx": \{
    "version": 1,
    "unlock_time": 60,
    "vin": [ \{
        "gen": \{
          "height": 0
        \}
      \}
    ],
    "vout": [ \{
        "amount": 17592186044415,
        "target": \{
          "key": "9b2e4c0281c0b02e7c53291a94d1d0cbff8883f8024f5142ee494ffbbd088071"
        \}
      \}
    ],
    "extra": [ 1, 119, 103, 170, 252, 222, 155, 224, 13, 207, 208, 152, 113, 94, 188, 
    247, 244, 16, 218, 235, 197, 130, 253, 166, 157, 36, 162, 142, 157, 11, 200, 144, 
    209
    ],
    "signatures": [ ]
  \},
  "tx_hashes": [ ]
\}
\end{Verbatim}



\section*{创世区块组件}

由于我们使用相同的软件打印创世区块和附录 \ref{appendix:block-content} 中的区块，结构看起来基本相同。我们指出一些独特的差异。

\begin{itemize}
	\item {\tt difficulty}（第 9 行）- 创世区块的难度报告为 1，这意味着任何 nonce 都可以。
	\item {\tt timestamp}（第 14 行）- 创世区块没有有意义的时间戳。
	\item {\tt prev\_id}（第 15 行）- 按惯例，我们使用空的 32 字节作为前一个区块 ID。
	\item {\tt nonce}（第 16 行）- 按惯例 $n = 10000$。
	\item {\tt amount}（第 27 行）- 这与我们在第 \ref{subsec:block-reward} 节中计算的第一个区块奖励（17.592186044415 XMR）完全一致。
	\item {\tt key}（第 29 行）- 最早的 Moneroj 被分配给 Monero 的创始人 thankful\_for\_today。
	\item {\tt extra}（第 33-36 行）- 使用附录 \ref{appendix:RCTTypeBulletproof2} 中讨论的编码方式，创世区块的矿工交易 {\tt extra} 字段只包含一个交易公钥。
	\item {\tt signatures}（第 37 行）- 创世区块中没有签名。这里是 {\tt print\_block} 函数的产物。第 39 行的 {\tt tx\_hashes} 也是如此。
\end{itemize}


\end{appendices}
